% AI-produced by GPT-6 Astra (OpenAI, ChatGPT Pro mode). No priority claim.
\subsection{Reduction of all actions to open stabilizers}

We first display the exact stabilizer hypothesis needed for arbitrary
actions. The preceding finite-support argument verifies that hypothesis,
and the final theorem combines the two proofs.

\begin{proposition}[Proved: an all-action stabilizer criterion]
\label{A:predimension-stabilizer-criterion}
Let $A$ be either ordered predimension category of
Proposition~\ref{A:predimension-quasi}. Let $M$ be its countable
strongly homogeneous rich limit, let $G=\operatorname{Aut}(M)$ with
the topology of pointwise convergence, and for a finite $d$-closed
substructure $B\subseteq M$ put
\[
 G_B=\{g\in G:g|_B=\id_B\}.
\]
Assume that every open subgroup of $G$ equals $G_F$ for some finite
$d$-closed $F\subseteq M$, allowing $F=\varnothing$.
Then $A$ is total. This conclusion concerns every covariant
$H:A\to\Set$, with no finiteness or injectivity assumption on $H$.
\end{proposition}

\begin{proof}
We give all changes of indexing category and all comparisons.
We use Quillen's Theorem A in its covariant-comma form: a functor
$u:C\to D$ with every $d\downarrow u$ weakly contractible induces
a weak equivalence on classifying spaces. This is also the immediate
contractible-fiber case of the opened
\citet[Theorem~IV.5.6]{GoerssJardine2009}: the comma transition maps
are then weak equivalences, and the displayed homotopy-cartesian
square has contractible fiber and contractible lower-left corner.

\emph{The rich limit and the concrete indexing category.}
There are countably many finite objects and arrows up to their finite
labels. The closure calculus, joint embedding, and amalgamation of
Lemma~\ref{A:predimension-amalgam} permit the usual rich-chain
construction: successively include each finite object and serve every
finite strong extension of every earlier stage, repeating the countable
enumeration. The union $M$ has the extension property for any finite
closed $B\subseteq M$: if $B\leq_d C$, first amalgamate $C$ with a
stage containing $B$, then use the served extension of that stage.
Thus $C$ strongly embeds into $M$ over $B$. The alternating
back-and-forth extension of finite closed partial isomorphisms shows
that every isomorphism between finite closed substructures of $M$
extends to an automorphism. This also follows from the opened strong
Fra\"iss\'e theorem, \citet[Theorem~2.6]{EvansHubickaNesetril2019};
the construction just given specifies its hypotheses here.

Let $A_M$ have as objects the nonempty finite $d$-closed
substructures of $M$, and as arrows \emph{all} increasing induced
$d$-closed embeddings between them, not merely inclusions. Every
finite object of $A$ is isomorphic to an object of $A_M$, and the
hom-sets are unchanged by this identification. Hence $A_M$ is
equivalent to $A$. Transport along this equivalence induces
equivalences of the corresponding categories of elements, so it
suffices to prove the assertion for $A_M$.

We also need the fixed-point identity
\[
                       \operatorname{Fix}_M(G_B)=B.
 \tag{F}\label{eq:predimension-stabilizer-fixed}
\]
If $x\notin B$, take a finite closed $D$ containing $B\cup\{x\}$.
Freely amalgamate two copies of $D$ over $B$, extend their common
order, and embed this finite strong extension into $M$ over the first
$D$. The two copies of $x$ are distinct. Their closed-copy
isomorphism extends to an automorphism fixing $B$, and moves $x$.
This proves \eqref{eq:predimension-stabilizer-fixed}, also for
$B=\varnothing$.

\emph{Nonempty cosieves in slices are contractible.}
Fix $P\in A_M$, and let $\mathcal U$ be a nonempty full cosieve
in $P/A_M$. Choose one of its objects $P\to Z$.
For $P\to X$ freely amalgamate $X$ and $Z$ over $P$.
In each gap of the order on $P$, put the $X$-side points before
the $Z$-side points, retaining the order within each side. On a
slice morphism $X\to Y$, use its given map on the $X$-side and
the identity on $Z$. This gives an increasing induced embedding.
It is $d$-closed because the target is the free amalgam of $Y$
and the source over $X$, whose image is closed in both. Thus this
construction is a functor on $P/A_M$, with natural maps from the
identity functor and from the constant object $P\to Z$.
Its image on $\mathcal U$ remains in $\mathcal U$, since it
receives a map from $P\to Z$. These natural transformations
contract $B\mathcal U$.

For precision about labels, this construction may first be made in
the concrete category of all finite tagged structures. Transporting
the result along an equivalence with $A_M$, using chosen
isomorphisms, gives a functor and natural transformations on the
slice in question. No functorial choice of an embedding into $M$
is being assumed. The same disjoint-union construction over the
empty structure, with one fixed nonempty object, contracts $BA_M$.

\emph{Removal of noninjective transitions.}
For an arbitrary $H:A_M\to\Set$ define an equivalence relation on
$H(a)$ by
\[
 h\sim_a h'\quad\Longleftrightarrow\quad
 \text{some }f:a\to b\text{ satisfies }H(f)h=H(f)h'.
\]
Reflexivity and symmetry are immediate. For transitivity, complete
two witnessing arrows with source $a$ to a commuting square by I.
Its two composites give one arrow witnessing transitivity. The
same square argument proves that $H(u)$ preserves this equivalence
relation for every $u:a\to a'$. Hence the quotients form a functor
$H^{\mathrm{inj}}$, with a natural surjection $q:H\to H^{\mathrm{inj}}$.
Every transition of $H^{\mathrm{inj}}$ is injective: a witness
after $u$ composes with $u$ to witness equality before it.

The induced functor
\[
             \El(q):\El_{A_M}H\longrightarrow\El_{A_M}H^{\mathrm{inj}}
\]
is homotopy final. To prove this, fix $(a,[h])$ in its target and
let $\mathcal D=(a,[h])\mathbin\downarrow\El(q)$.
There is a functor $J:a/A_M\to\mathcal D$ sending
$f:a\to x$ to $((x,H(f)h),f)$. For
$d=((x,k),f)\in\mathcal D$, its comma category $d\downarrow J$
is exactly the full cosieve in $x/A_M$ consisting of arrows
$g:x\to y$ satisfying
\[
                         H(g)k=H(gf)h.
\]
It is nonempty by $[k]=[H(f)h]$. It is contractible by the preceding
slice-cosieve construction. Quillen's Theorem A therefore makes
$BJ$ a weak equivalence. Since $a/A_M$ has an initial object,
$B\mathcal D$ is weakly contractible. A second application of
Theorem A proves that $B\El(q)$ is a weak equivalence. It thus
suffices to consider actions with injective transitions.

\emph{Saturation of an injective action.}
Let $H$ now have injective transitions. The finite closed
substructures of $M$, with inclusion maps only, form a directed
poset $\mathcal I$: every finite union is contained in one of the
finite closed stages. Put
\[
                         S=\operatorname{colim}_{B\in\mathcal I} H(B).
\]
All canonical maps $i_B:H(B)\to S$ are injective. An automorphism
$g\in G$ acts by
\[
                   g\cdot i_B(h)=i_{gB}\bigl(H(g|_B)h\bigr).
\]
Naturality with respect to inclusions makes this well-defined, and
functoriality gives a group action. Every $i_B(h)$ is fixed by
$G_B$, so $S$ is a continuous discrete $G$-set.

For any such $G$-set $S$, define
\[
                         T_S(B)=S^{G_B}.
\]
For an arrow $f:B\to C$, extend $f$ to $g\in G$ and set
$T_S(f)(s)=g\cdot s$. This is independent of the extension:
two extensions differ on the right by an element of $G_B$.
Its value is fixed by $G_C$, since $G_C\subseteq G_{fB}$.
Extending composites proves functoriality. The maps $i_B$ above
give a natural monomorphism $i:H\hookrightarrow T_S$.

For $(a,s)\in\El T_S$, the comma category
$(a,s)\downarrow\El(i)$ identifies with the full cosieve of
$a/A_M$ on those $f:a\to b$ for which
$T_S(f)(s)$ belongs to $i_b(H(b))$. The corresponding element of
$H(b)$ is unique by injectivity of $i_b$, so this identification
also preserves all arrows. The cosieve is nonempty: represent
$s$ by some $h\in H(c)$ and choose a finite closed $b$ containing
$a\cup c$. For the inclusion $a\hookrightarrow b$, the identity
automorphism of $M$ is an extension, and $s$ is represented in
$H(b)$ by the inclusion from $c$. The slice-cosieve contraction
and Theorem A now give a weak equivalence
\[
                      B\El H\longrightarrow B\El T_S.
\]

\emph{The open-subgroup hypothesis finishes the argument.}
Decompose $S$ into its $G$-orbits. Since every stabilizer is open,
the stated hypothesis identifies each orbit with $G/G_F$ for a
finite closed $F$, possibly empty. Taking fixed subsets preserves
this orbit decomposition. If $F\ne\varnothing$, there is a natural
isomorphism of $A_M$-sets
\[
                   T_{G/G_F}(B)\ \cong\ A_M(F,B).
\]
Indeed, a coset $gG_F$ is fixed by $G_B$ precisely when every
element of $G_B$ fixes $gF$ pointwise. By
\eqref{eq:predimension-stabilizer-fixed}, this is equivalent to
$gF\subseteq B$. The map $g|_F:F\to B$ is a closed embedding by
restriction of closure, every such embedding extends to $G$, and
two extensions give the same coset exactly when their restrictions
to $F$ agree. Postcomposition agrees with the definition of $T_S$.
The category of elements of this representable has the initial
object $(F,\id_F)$. If $F=\varnothing$, the orbit is a singleton,
$T_{G/G_F}$ is constant one, and its category of elements is the
already contractible category $A_M$.

Thus every component of $B\El T_S$ is contractible. The two
proved weak equivalences transfer this assertion first to every
injective $H$ and then to every $H$. Nonemptiness and connectedness
of $A_M$ follow already from joint embedding. This is totality.
\end{proof}

\begin{theorem}[Proved: a total category with an all-ring triple obstruction]
\label{A:predimension-total}
Both $\mathcal C_\delta$ and its full subcategory whose raw points
and hyperedges are all $d$-closed are total. For each of these
categories and every covariant $H:A\to\Set$, every component of
$\El_AH$ is weakly contractible. In each category the three point
arrows $V\rightrightarrows W$ of
Proposition~\ref{A:predimension-triple} have no affine equalizer
over any nonzero unital ring $R$, and the displayed column cokernel
$Q_R$ satisfies $T_1(Q_R)\cong R$. Thus both categories are
separators, simultaneously for integral, rational and every prime-field
coefficient system.
\end{theorem}
\begin{proof}
The closure and amalgamation lemma constructs the required rich
homogeneous limit for each of the two classes. The alternating
bridge of Lemma~\ref{A:predimension-alternating-bridge} proves
the stabilizer-generation identity of
Proposition~\ref{A:predimension-stabilizers}. Minimal finite
supports and rigidity of finite linear orders then give
Proposition~\ref{A:predimension-open-subgroups}: every open
subgroup is the pointwise stabilizer of a finite closed subset,
including the empty subset. This is precisely the hypothesis of
Proposition~\ref{A:predimension-stabilizer-criterion}, whose two
homotopy-final comparisons apply to all actions. It proves totality.
The objects $V,W$ belong to both classes, their three point maps
are closed embeddings, and
Proposition~\ref{A:predimension-triple} proves the remaining
assertions over every stated coefficient ring. No restriction to
finite actions or to injective transition maps remains.
\end{proof}
